% =====================================================================
%  Bab II — Tinjauan Pustaka (rekonstruksi dari naskah PDF).
%  Pemeriksaan ulang oleh penulis dianjurkan (transkripsi dari PDF).
% =====================================================================

\section{Tinjauan Pustaka}

\subsection{Internet of Things (IoT) dan Ekosistem Smart Home}
Konsep Internet of Things (IoT) pertama kali dipopulerkan oleh Kevin Ashton
pada tahun 1999 dalam konteks manajemen rantai pasok. Definisi yang lebih
komprehensif dikemukakan oleh Atzori \textit{et al.}~\cite{b5}, yang
mendeskripsikan IoT sebagai sebuah paradigma teknologi di mana objek-objek
fisik di dunia nyata --- yang dilengkapi dengan sensor, aktuator, dan
kapabilitas komunikasi --- saling terhubung melalui jaringan internet untuk
mengumpulkan, bertukar, dan memproses data secara otonom.

Dalam konteks rumah tangga, IoT diwujudkan melalui ekosistem smart home yang
mengintegrasikan berbagai perangkat cerdas ke dalam satu jaringan domestik.
Al-Fuqaha \textit{et al.}~\cite{b6} mengklasifikasikan arsitektur komunikasi
IoT ke dalam beberapa lapisan: lapisan persepsi (sensor dan aktuator); lapisan
jaringan yang menangani transmisi data melalui protokol seperti Wi-Fi (IEEE
802.11), Zigbee (IEEE 802.15.4), Z-Wave, BLE, dan LoRaWAN; lapisan middleware;
serta lapisan aplikasi. Berdasarkan data ITU~\cite{b7}, penetrasi perangkat IoT
di rumah tangga meningkat rata-rata 23\% per tahun dalam periode 2020--2024.
Setiap perangkat berjalan dengan sistem operasi dan firmware yang bervariasi,
menggunakan protokol komunikasi yang heterogen, dan memiliki tingkat keamanan
yang tidak seragam --- menciptakan permukaan serangan yang luas dan kompleks.

\subsection{Ancaman dan Kerentanan Keamanan pada Perangkat IoT}

\subsubsection{Taksonomi Ancaman Keamanan IoT}
Neshenko \textit{et al.}~\cite{b8} menyajikan taksonomi ancaman keamanan IoT
yang komprehensif melalui survei terhadap lebih dari 200 literatur akademik,
mengkategorikan ancaman ke dalam empat lapisan utama:
\begin{itemize}
    \item \textbf{Lapisan Perangkat:} kerentanan pada firmware, hardware
    (debug interface seperti JTAG dan UART), serta konfigurasi bawaan pabrik
    yang tidak aman.
    \item \textbf{Lapisan Jaringan:} serangan Man-in-the-Middle (MitM), packet
    sniffing, denial of service (DoS), dan eksploitasi terhadap port dan
    layanan jaringan yang rentan.
    \item \textbf{Lapisan Cloud:} kerentanan pada backend services, API yang
    tidak aman, dan penyimpanan data yang tidak terenkripsi.
    \item \textbf{Lapisan Kecerdasan Buatan:} ancaman terhadap sistem deteksi
    berbasis machine learning seperti adversarial attacks dan model poisoning.
\end{itemize}
Temuan empiris mereka menunjukkan bahwa 78\% eksploitasi IoT berskala internet
menargetkan kerentanan pada lapisan jaringan, menjadikan analisis keamanan pada
tingkat jaringan sebagai prioritas utama dalam perlindungan perangkat IoT
domestik.

\subsubsection{Kerangka Acuan OWASP IoT Top 10}
Open Web Application Security Project (OWASP) menerbitkan IoT Top 10~\cite{b9}
sebagai kerangka acuan yang mengidentifikasi sepuluh kategori kerentanan paling
kritis pada ekosistem IoT:
\begin{enumerate}
    \item \textit{Weak, Guessable, or Hardcoded Passwords:} penggunaan kata
    sandi lemah, mudah ditebak, atau hardcoded dalam firmware.
    \item \textit{Insecure Network Services:} layanan jaringan yang tidak
    diperlukan atau tidak dikonfigurasi dengan aman, yang terekspos.
    \item \textit{Insecure Ecosystem Interfaces:} kerentanan pada antarmuka
    web, API, mobile app, atau cloud backend.
    \item \textit{Lack of Secure Update Mechanism:} ketiadaan mekanisme
    pembaruan firmware yang aman.
    \item \textit{Use of Insecure or Outdated Components:} penggunaan komponen
    perangkat lunak pihak ketiga yang memiliki kerentanan yang telah diketahui.
    \item \textit{Insufficient Privacy Protection:} perlindungan privasi yang
    tidak memadai terhadap data pribadi pengguna.
    \item \textit{Insecure Data Transfer and Storage:} transfer dan penyimpanan
    data yang tidak terenkripsi atau menggunakan enkripsi yang lemah.
    \item \textit{Lack of Device Management:} ketiadaan mekanisme pengelolaan
    perangkat yang memadai pasca-deployment.
    \item \textit{Insecure Default Settings:} konfigurasi bawaan pabrik yang
    tidak aman, termasuk default password dan layanan yang aktif otomatis.
    \item \textit{Lack of Physical Hardening:} kurangnya perlindungan terhadap
    akses fisik yang tidak sah ke antarmuka hardware perangkat.
\end{enumerate}
Dalam penelitian ini, OWASP IoT Top 10 digunakan sebagai basis pemetaan
kerentanan yang terdeteksi oleh aplikasi pemindai terhadap kategori-kategori
risiko yang terstandarisasi.

Meskipun diterbitkan pada tahun 2018, OWASP IoT Top 10 tetap menjadi kerangka
acuan yang paling banyak dikutip dalam literatur akademik keamanan IoT hingga
saat ini. Hal ini disebabkan oleh dua faktor: pertama, OWASP Foundation belum
menerbitkan versi pembaruan resmi untuk kategori IoT --- berbeda dengan OWASP
Web Top 10 dan Mobile Top 10 yang telah diperbarui secara berkala; kedua,
kategori kerentanan yang diidentifikasi pada versi 2018 --- seperti penggunaan
kredensial bawaan (I1), layanan jaringan tidak aman (I2), dan transfer data
tidak terenkripsi (I7) --- terbukti tetap relevan karena masih menjadi vektor
eksploitasi utama pada insiden keamanan IoT kontemporer, termasuk varian botnet
pasca-Mirai yang terus mengeksploitasi pola kerentanan yang sama.

\subsubsection{Studi Kasus: Mirai Botnet}
Antonakakis \textit{et al.}~\cite{b1} menyajikan analisis teknis yang
mengungkapkan mekanisme infeksi Mirai: botnet ini secara sistematis memindai
ruang alamat IPv4 untuk menemukan perangkat IoT dengan layanan Telnet (port 23
dan 2323) yang terbuka, kemudian mencoba masuk menggunakan daftar 62 kombinasi
username dan password bawaan pabrik yang di-hardcode. Kolias
\textit{et al.}~\cite{b10} mengelaborasi bahwa serangan DDoS terhadap penyedia
DNS Dyn pada Oktober 2016 menghasilkan volume lalu lintas hingga 1,2 Tbps.
Meidan \textit{et al.}~\cite{b22} mengusulkan pendekatan deteksi botnet IoT
berbasis deep autoencoder (N-BaIoT) yang mendeteksi aktivitas Mirai dan
BASHLITE dengan akurasi di atas 99\%.

\subsubsection{Kerentanan Protokol Wi-Fi}
Vanhoef dan Piessens~\cite{b11} menemukan kerentanan fundamental pada protokol
WPA2 yang dinamakan Key Reinstallation Attack (KRACK), yang mengeksploitasi
kelemahan dalam mekanisme four-way handshake. Sebagai respons, Wi-Fi Alliance
merilis WPA3; namun Vanhoef dan Ronen~\cite{b12} dalam penelitian Dragonblood
mengungkapkan bahwa implementasi awal handshake Dragonfly pada WPA3 rentan
terhadap serangan side-channel dan downgrade.

\subsection{Teknik Pemindaian dan Deteksi Kerentanan Jaringan}

\subsubsection{Prinsip Network Scanning}
Lyon~\cite{b13} menjelaskan bahwa pemindaian jaringan umumnya terdiri atas:
(a) \textit{Host Discovery} --- identifikasi host aktif, umumnya menggunakan
ARP scanning, ICMP echo request, atau TCP SYN/ACK probing; (b) \textit{Port
Scanning} --- pemeriksaan status port melalui TCP SYN scan, TCP connect scan,
dan UDP scan; (c) \textit{Service Detection} --- identifikasi layanan dan versi
melalui banner grabbing; serta (d) \textit{OS Detection}. Nmap, yang
dikembangkan oleh Gordon Lyon, merupakan tools pemindai jaringan yang paling
banyak digunakan~\cite{b13}, dan dalam penelitian ini Nmap diintegrasikan
secara opsional untuk memperkaya hasil pemindaian.

\subsubsection{Device Fingerprinting pada Perangkat IoT}
Sivanathan \textit{et al.}~\cite{b14} mengembangkan metode klasifikasi
perangkat IoT berbasis karakteristik lalu lintas jaringan dan berhasil
mengklasifikasikan 28 jenis perangkat dengan akurasi mencapai 99,88\%
menggunakan Random Forest. Dalam konteks aplikasi pemindai keamanan rumah,
teknik fingerprinting disederhanakan menjadi identifikasi berbasis MAC address
(vendor lookup via OUI) dan karakteristik layanan yang terdeteksi (mis. RTSP
pada port 554 mengindikasikan kamera IP, UPnP mengindikasikan smart device).

\subsubsection{Metodologi Vulnerability Assessment}
Weidman~\cite{b15} membedakan pendekatan aktif dan pasif dalam penilaian
kerentanan. Dalam konteks pemindaian jaringan rumah, pendekatan semi-aktif ---
mencakup pemindaian port dan deteksi layanan tanpa eksploitasi --- merupakan
kompromi yang tepat antara kedalaman analisis dan keselamatan operasional.
Shah dan Issac~\cite{b16} menyimpulkan bahwa integrasi pendekatan berbasis
aturan (rule-based) dengan machine learning menghasilkan akurasi deteksi yang
optimal; temuan ini menginformasikan kombinasi aturan pemetaan berbasis OWASP
IoT Top 10 dengan logika skoring risiko berbasis heuristik dalam penelitian
ini.

\subsection{Kerangka Kerja Keamanan IoT}

\subsubsection{NIST Foundational Cybersecurity Activities for IoT}
NIST melalui dokumen NISTIR 8259~\cite{b17} menetapkan enam aktivitas
fundamental keamanan siber untuk perangkat IoT. HomeGuard difokuskan untuk
memfasilitasi dua aktivitas yang paling relevan bagi pengguna rumahan, yaitu
identifikasi aset (melalui network discovery dan device profiling) dan deteksi
kerentanan (melalui port scanning, service detection, dan pemetaan OWASP IoT
Top 10).

\subsubsection{Defense-in-Depth untuk Jaringan Rumah}
Sicari \textit{et al.}~\cite{b18} mengemukakan bahwa keamanan IoT harus
diterapkan melalui pendekatan berlapis (defense-in-depth). Dalam kerangka ini,
aplikasi pemindai keamanan jaringan berperan sebagai komponen ``deteksi'' pada
lapisan jaringan, yang memberikan visibilitas kepada pengguna mengenai kondisi
keamanan jaringan dan perangkat IoT yang terhubung.

\subsubsection{Arsitektur Zero Trust untuk IoT}
Rose \textit{et al.}~\cite{b19} memperkenalkan konsep Zero Trust Architecture
(ZTA) dengan prinsip ``never trust, always verify''. Penerapannya pada jaringan
rumah mengimplikasikan bahwa setiap perangkat IoT harus diperlakukan sebagai
entitas yang berpotensi tidak tepercaya, sehingga visibilitas menyeluruh ---
yang difasilitasi oleh tools pemindai --- menjadi prasyarat fundamental.

\subsection{Penelitian Terdahulu yang Relevan}
Beberapa penelitian terdahulu telah mengeksplorasi keamanan IoT rumah tangga.
Alrawi \textit{et al.}~\cite{b2} melakukan evaluasi sistematis (SoK) namun tidak
menyertakan tools automasi untuk pengguna akhir. Hafeez
\textit{et al.}~\cite{b4} mengembangkan IoT-KEEPER untuk continuous monitoring
di edge, yang kurang praktis untuk pemindaian ad-hoc. Acar
\textit{et al.}~\cite{b20} (Peek-a-Boo) menyoroti risiko privasi pada lalu
lintas terenkripsi. Hamza \textit{et al.}~\cite{b21} mengusulkan deteksi
berbasis MUD + SDN yang bergantung pada ketersediaan profil MUD. Tabel~\ref{tab:terdahulu}
menyajikan perbandingan posisi penelitian ini.

\begin{table*}[t]
    \centering
    \caption{Perbandingan Penelitian Terdahulu}
    \label{tab:terdahulu}
    \footnotesize
    \begin{tabular}{|c|l|c|l|l|l|}
        \hline
        \textbf{No} & \textbf{Peneliti} & \textbf{Tahun} & \textbf{Metode} &
        \textbf{Fokus Utama} & \textbf{Perbedaan dengan Penelitian Ini} \\
        \hline
        1 & Antonakakis \textit{et al.}~\cite{b1} & 2017 & Analisis forensik botnet & Mekanisme infeksi Mirai & Fokus post-incident, bukan pencegahan \\
        \hline
        2 & Alrawi \textit{et al.}~\cite{b2} & 2019 & Evaluasi sistematis (SoK) & Keamanan IoT rumah (4 dimensi) & Survey/taxonomy, tanpa tools automasi \\
        \hline
        3 & Hafeez \textit{et al.}~\cite{b4} & 2020 & Analisis lalu lintas online & Deteksi anomali di edge & Continuous monitoring, bukan ad-hoc \\
        \hline
        4 & Sivanathan \textit{et al.}~\cite{b14} & 2019 & ML classification & Klasifikasi perangkat IoT & Fokus identifikasi, bukan kerentanan \\
        \hline
        5 & Acar \textit{et al.}~\cite{b20} & 2020 & Analisis traffic terenkripsi & Inferensi aktivitas smart home & Fokus privasi, bukan vulnerability scanning \\
        \hline
        6 & Hamza \textit{et al.}~\cite{b21} & 2020 & SDN + MUD profiling & Deteksi serangan volumetrik & Bergantung pada profil MUD produsen \\
        \hline
        7 & Meidan \textit{et al.}~\cite{b22} & 2018 & Deep autoencoder & Deteksi botnet IoT & ML-based, butuh data pelatihan \\
        \hline
        \textbf{8} & \textbf{Penelitian ini} & \textbf{2025} & Python + Nmap + OWASP & End-user scanner & Tools praktis berbasis GUI dengan pemetaan OWASP IoT Top 10 \\
        \hline
    \end{tabular}
\end{table*}
